% Part: turing-machines
% Chapter: machines-computations
% Section: representing-tms

\documentclass[../../../include/open-logic-section]{subfiles}

\begin{document}

% SEGMENT OLP-0255-S01

\olfileid{tur}{mac}{rep}
\olsection{டியூரிங் பொறிகளைக் குறிப்பிடுதல்}

\begin{explain}
டியூரிங் பொறிகளை \emph{நிலை வரைபடங்களால்} காட்சிப்படுத்தலாம்.
அம்புகளால் இணைக்கப்பட்ட நிலைக் கலங்கள் இவ்வரைபடங்களை உருவாக்குகின்றன.
எதிர்பார்க்கத்தக்கவாறே, ஒவ்வொரு நிலைக் கலமும் பொறியின் ஒரு நிலையைக்
குறிக்கிறது. ஒவ்வொரு அம்பும் அந்நிலையிலிருந்து நிறைவேற்றக்கூடிய ஒரு
கட்டளையைக் குறிக்கிறது; கட்டளையின் விவரங்கள் அதற்குரிய அம்பின்
மேலோ கீழோ எழுதப்படுகின்றன. $q_0$, $q_1$ என்ற இரு உள்நிலைகளையும்
ஒரே ஒரு கட்டளையையும் மட்டும் கொண்ட பின்வரும் பொறியைக் கருதுக:
\[
\begin{tikzpicture}[->,>=stealth',shorten >=1pt,auto,node distance=2.8cm,
                    semithick]
  \tikzstyle{every state}=[fill=none,draw=black,text=black]

  \node[initial,state]   (A)              {$q_0$};
  \node[state]   (B) [right of=A] {$q_1$};

  \path (A) edge  node {\TMtrans{\TMblank}{\TMstroke}{\TMright}} (B);
\end{tikzpicture}
\]
டியூரிங் பொறிக்கு ஒரு வாசி-எழுது தலையும், அதன்மீது உள்ளீடு எழுதப்பட்ட
ஒரு நாடாவும் உள்ளன என்பதை நினைவுகூர்க. இக்கட்டளையைப் பின்வருமாறு
வாசிக்கலாம்: \emph{நிலை $q_0$ இல் ஒரு~$\TMblank$ ஐ வாசித்தால்,
ஒரு~$\TMstroke$ ஐ எழுதி, வலப்புறம் நகர்ந்து, நிலை $q_1$ குச் செல்க}.
இது நிலைமாற்றுச் சார்பு $\tuple{q_0, \TMblank}$ ஐ
$\tuple{q_1, \TMstroke, \TMright}$ கு அனுப்புவதற்குச் சமானம்.
\end{explain}

\begin{ex}
\emph{இரட்டை எண் பொறி}: நாடாவில் முதல் $\TMblank$ க்கு முன்
எல்லா $\TMstroke$-களும் வருகின்றன என்ற கருதுகோளின்கீழ்,
நாடாவில் இரட்டை எண்ணிக்கையிலான $\TMstroke$-கள் இருந்தால், எனில்
மட்டுமே, பின்வரும் டியூரிங் பொறி நிற்கிறது.
\[
\begin{tikzpicture}[->,>=stealth',shorten >=1pt,auto,node distance=2.8cm,
                    semithick]
  \tikzstyle{every state}=[fill=none,draw=black,text=black]

  \node[initial,state]         (A)              {$q_0$};
  \node[state]         (B) [right of=A] {$q_1$};

  \path (A) edge [bend left] node {\TMtrans{\TMstroke}{\TMstroke}{\TMright}} (B)
        (B) edge [loop above] node {\TMtrans{\TMblank}{\TMblank}{\TMright}} (B)
            edge [bend left] node {\TMtrans{\TMstroke}{\TMstroke}{\TMright}} (A);
\end{tikzpicture}
\]

நிலை வரைபடம் பின்வரும் நிலைமாற்றுச் சார்புக்குரியது:
\begin{align*}
\delta(q_0, \TMstroke) & = \tuple{q_1, \TMstroke, \TMright},\\
\delta(q_1, \TMstroke) & = \tuple{q_0, \TMstroke, \TMright},\\
\delta(q_1, \TMblank)  & = \tuple{q_1, \TMblank, \TMright}
\end{align*}
\end{ex}

\begin{explain}
உள்ளீடு இரட்டை எண்ணிக்கையிலான கோடுகளாக இருக்கும்போது மட்டுமே
மேலுள்ள பொறி நிற்கிறது. இல்லையெனில், பொறி (கொள்கையளவில்)
முடிவின்றிச் செயல்படத் தொடர்கிறது. எந்தப் பொறிக்கும் உள்ளீட்டுக்கும்,
வெளியீட்டைத் தீர்மானிக்கப் பொறியின் \emph{அமைவுநிலைகளை}
ஒன்றன்பின் ஒன்றாகத் தடமறியலாம். அமைவுநிலைகளுக்கான முறைசார்ந்த
வரையறையைப் பின்னர் தருவோம். இப்போதைக்கு, பொறி இயங்கும்போது எந்த
நேரத்திலும் அதன் நிலையைப் படம்பிடிக்கும் வரைபடங்களின் தொடராக
அமைவுநிலைகளை உள்ளுணர்வாகக் கருதலாம். நாடாவின் உள்ளடக்கம்,
பொறியின் நிலை, வாசி-எழுது தலையின் இருப்பிடம் ஆகியவற்றை
அமைவுநிலைகள் காட்டுகின்றன.

நான்கு $\TMstroke$-களை உள்ளீடாகக் கொண்டு தொடங்கும்போது, இரட்டை
எண் பொறியின் அமைவுநிலைகளை ஒன்றன்பின் ஒன்றாகத் தடமறிவோம்.
இந்நேர்வில் பொறி நிற்கும் என்று எதிர்பார்க்கிறோம். பிறகு மூன்று
$\TMstroke$-களை உள்ளீடாகக் கொண்டு பொறியை இயக்குவோம்; அங்கு பொறி
என்றென்றும் இயங்கும்.

பொறி நிலை~$q_0$ இல் தொடங்கி, இடக்கோடி~$\TMstroke$ ஐ வருடுகிறது.
பொறியின் தொடக்க நிலையைப் பின்வருமாறு குறிப்பிடலாம்:
\[
\TMendtape \TMstroke_0 \TMstroke \TMstroke \TMstroke \TMblank \ldots
\]
மேலுள்ள அமைவுநிலை நேரடியானது. காணக்கூடியவாறு, பொறி நிலை ஒன்றில்
தொடங்கி, இடக்கோடி~$\TMstroke$ ஐ வருடுகிறது. முதல்~$\TMstroke$ இல்
நிலைப் பெயரை கீழொட்டாகக் குறிப்பதால் இது காட்டப்படுகிறது.
இப்புள்ளியில் பொருந்தும் கட்டளை
$\delta(q_0, \TMstroke) = \tuple{q_1, \TMstroke, \TMright}$;
எனவே பொறி நாடாவில் வலப்புறம் நகர்ந்து நிலை~$q_1$ கு மாறுகிறது.
\[
\TMendtape \TMstroke \TMstroke_1 \TMstroke \TMstroke \TMblank \ldots
\]
பொறி இப்போது நிலை~$q_1$ இல் ஒரு~$\TMstroke$ ஐ வருடுவதால்,
$\delta(q_1, \TMstroke) = \tuple{q_0, \TMstroke, \TMright}$
என்ற கட்டளையை நாம் ``பின்பற்ற'' வேண்டும். இதனால்
பின்வரும் அமைவுநிலை கிடைக்கிறது:
\[
\TMendtape \TMstroke \TMstroke \TMstroke_0 \TMstroke \TMblank \ldots
\]
பொறி தொடரும்போது, விதிகள் மீண்டும் அதே வரிசையில் பயன்படுத்தப்பட்டு,
பின்வரும் இரு அமைவுநிலைகள் கிடைக்கின்றன:
\[
\TMendtape \TMstroke \TMstroke \TMstroke \TMstroke_1 \TMblank \ldots
\]
\[
\TMendtape \TMstroke \TMstroke \TMstroke \TMstroke \TMblank_0 \ldots
\]
பொறி இப்போது நிலை~$q_0$ இல் ஒரு~$\TMblank$ ஐ வருடுகிறது.
நிறைவேற்ற வேண்டிய கட்டளை எதுவும் இல்லை என்பதை நிலைமாற்று
வரைபடத்திலிருந்து எளிதில் காணலாம்; எனவே பொறி நின்றுள்ளது.
உள்ளீடு ஏற்கப்பட்டது என்பதே இதன் பொருள்.

அடுத்து, மூன்று $\TMstroke$-களை உள்ளீடாகக் கொண்டு பொறியைத்
தொடங்குகிறோம் எனக் கொள்க. அதே கட்டளைகள் நிறைவேற்றப்படுவதால்
முதல் சில அமைவுநிலைகள் ஒத்தவை; நாடா உள்ளீட்டில் மட்டும் சிறு
வேறுபாடு உள்ளது:
\[
\TMendtape \TMstroke_0 \TMstroke \TMstroke \TMblank \ldots
\]
\[
\TMendtape \TMstroke \TMstroke_1 \TMstroke \TMblank \ldots
\]
\[
\TMendtape \TMstroke \TMstroke \TMstroke_0 \TMblank \ldots
\]
\[
\TMendtape \TMstroke \TMstroke \TMstroke \TMblank_1 \ldots
\]
பொறி இப்போது எல்லா $\TMstroke$-களையும் கடந்து, நிலை~$q_1$ இல்
ஒரு~$\TMblank$ ஐ வாசிக்கிறது. வரைபடத்தில் காட்டப்பட்டவாறு,
$\delta(q_1, \TMblank) =\tuple{q_1, \TMblank, \TMright}$ என்ற
வடிவில் ஒரு கட்டளை உள்ளது. நாடா வலப்புறத்தில் முடிவின்றி
$\TMblank$ ஆல் நிரப்பப்பட்டுள்ளதால், பொறி இக்கட்டளையை
\emph{என்றென்றும்} நிறைவேற்றத் தொடரும்; நிலை~$q_1$ இலேயே இருந்து
மேன்மேலும் வலப்புறம் நகரும். பொறி ஒருபோதும் நிற்காது; உள்ளீட்டை
ஏற்காது.
\end{explain}

\begin{explain}
எல்லாப் பொறிகளும் நிற்காது என்பதைக் கவனிப்பது முக்கியம். நிற்றல்
என்பது, நிறைவேற்ற வேண்டிய கட்டளைகள் பொறியிடம் தீர்ந்துவிடுவது
என்றால், ஒவ்வொரு நிலையிலும் ஒவ்வொரு குறிக்கும் வெளியேறும் அம்பு
இருப்பதை உறுதிசெய்து ஒருபோதும் நிற்காத பொறியை உருவாக்கலாம்.
~$q_0$ இல் ஒரு $\TMblank$ ஐ வருடுவதற்கான கட்டளையைச் சேர்த்து,
இரட்டை எண் பொறியை முடிவின்றி இயங்குமாறு மாற்றலாம்.
\end{explain}

\begin{ex}
\[
\begin{tikzpicture}[->,>=stealth',shorten >=1pt,auto,node distance=2.8cm,
                    semithick]
  \tikzstyle{every state}=[fill=none,draw=black,text=black]

  \node[initial,state] (A)              {$q_0$};
  \node[state]         (B) [right of=A] {$q_1$};

  \path
  (A) edge [bend left] node {\TMtrans{\TMstroke}{\TMstroke}{\TMright}} (B)
      edge [loop above] node {\TMtrans{\TMblank}{\TMblank}{\TMright}} (A)
  (B) edge [loop above] node {\TMtrans{\TMblank}{\TMblank}{\TMright}} (B)
      edge [bend left] node {\TMtrans{\TMstroke}{\TMstroke}{\TMright}} (A);
\end{tikzpicture}
\]
\end{ex}

\begin{explain}
டியூரிங் பொறிகளைக் குறிப்பிடுவதற்கு பொறி அட்டவணைகள்
மற்றொரு வழி. பொறி அட்டவணையில் நாடாவின் குறியெழுத்துத் தொகுதி
$x$-அச்சிலும், பொறி நிலைகளின் கணம் $y$-அச்சிலும் காட்டப்படுகின்றன.
அட்டவணைக்குள், ஒவ்வொரு நிலையும் குறியும் வெட்டும் இடத்தில்
கட்டளையின் மீதிப்பகுதி---புதிய நிலை, புதிய குறி, நகர்வுத் திசை---
எழுதப்படுகிறது. பொறி எந்த நிலையில், எந்தக் குறிக்காக நிற்கிறது
என்று தீர்மானிப்பதைப் பொறி அட்டவணைகள் எளிதாக்குகின்றன.
அட்டவணையில் ஓர் இடைவெளி இருக்கும் போதெல்லாம், அது பொறி
நிற்கக்கூடிய ஒரு புள்ளி. பொறி நிற்கும் புள்ளிகள் எப்போதும் உடனடியாகத்
தெளிவாக இல்லாத நிலை வரைபடங்கள், கட்டளைத் தொகுதிகள் ஆகியவற்றைப்
போலன்றி, பொறி அட்டவணையின் இடைவெளிகளைக் கண்டறிந்து எந்த நிற்றல்
புள்ளியையும் விரைவாக அடையாளம் காணலாம்.
\end{explain}

\begin{ex}
இரட்டை எண் பொறிக்கான பொறி அட்டவணை:
\[
\centering
\begin{tabular}{lllll}
\cline{2-4}
\multicolumn{1}{l|}{}      & \multicolumn{1}{c|}{$\TMblank$}
& \multicolumn{1}{c|}{$\TMstroke$}     & \multicolumn{1}{c|}{$\TMendtape$}   &  \\ \cline{1-4}
\multicolumn{1}{|l|}{$q_0$} & \multicolumn{1}{l|}{}
& \multicolumn{1}{l|}{\TMtrans{\TMstroke}{q_1}{\TMright}} &
\multicolumn{1}{l|}{\phantom{\TMtrans{\TMstroke}{q_1}{\TMright}}} &  \\ \cline{1-4}
\multicolumn{1}{|l|}{$q_1$} & \multicolumn{1}{l|}{\TMtrans{\TMblank}{q_1}{\TMright}}
& \multicolumn{1}{l|}{\TMtrans{\TMstroke}{q_0}{\TMright}} &
\multicolumn{1}{l|}{\phantom{\TMtrans{\TMstroke}{q_1}{\TMright}}} &  \\ \cline{1-4}
\end{tabular}
\]
நாம் காணக்கூடியவாறு, நிலை~$q_0$ இல் ஒரு~$\TMblank$ ஐ வருடும்போது
பொறி நிற்கிறது.
\end{ex}

\begin{explain}
இதுவரை உள்ளீட்டை வாசித்து ஏற்கும் பொறிகளை மட்டுமே கருதியுள்ளோம்.
ஆனால் டியூரிங் பொறிகளால் வாசிக்கவும் எழுதவும் முடியும்.
இத்தகைய பொறிக்கு ஓர் எடுத்துக்காட்டு (இன்னும் மிகப் பல
எடுத்துக்காட்டுகள் இருந்தாலும்) ஒரு \emph{இரட்டிப்பான்}.
நாடாவில்~$n$ $\TMstroke$-களின் ஒரு தொகுதியுடன் தொடங்கும்போது,
இரட்டிப்பான்~$2n$ $\TMstroke$-களின் ஒரு தொகுதியை வெளியிடுகிறது.
\end{explain}

\begin{ex}
  \ollabel{ex:doubler} இரட்டிப்பான் பொறியைக் கட்டுவதற்கு முன்,
  சிக்கலைத் தீர்ப்பதற்கான ஒரு \emph{செயலுத்தியை} வகுப்பது முக்கியம்.
  பொறி (நாம் அதை வகுத்துள்ளபடி) தான் எத்தனை $\TMstroke$-களை
  வாசித்துள்ளது என்பதை நினைவில் வைத்துக்கொள்ள முடியாது; எனவே
  நாடாவிலுள்ள எல்லா $\TMstroke$-களையும் கண்காணிக்க ஒரு வழியை
  வகுக்க வேண்டும். வெளியீட்டை உள்ளீட்டிலிருந்து ஒரு~$\TMblank$ ஆல்
  பிரிப்பது அத்தகைய ஒரு வழி. அப்போது பொறி உள்ளீட்டிலுள்ள முதல்
  $\TMstroke$ ஐ அழித்து, உள்ளீட்டின் மீதிப்பகுதியைக் கடந்து,
  ஒரு $\TMblank$ ஐ விட்டுவிட்டு, இரு புதிய $\TMstroke$-களை எழுதலாம்.
  பிறகு பொறி பின்னே சென்று, உள்ளீட்டிலுள்ள இரண்டாவது $\TMstroke$ ஐக்
  கண்டுபிடித்து அதையும் இரட்டிக்கும். உள்ளீட்டின் ஒவ்வொரு
  $\TMstroke$ க்கும், வெளியீட்டின் இரு $\TMstroke$-களை எழுதும்.
  பொறி செல்லச்செல்ல உள்ளீட்டை அழிப்பதால், எந்த $\TMstroke$ உம்
  தவறவிடப்படவோ இருமுறை இரட்டிக்கப்படவோ இல்லை என்று உறுதிசெய்யலாம்.
  முழு உள்ளீடும் அழிக்கப்பட்டதும், நாடாவில் $2n$ $\TMstroke$-கள்
  மீதமிருக்கும். இதனால் கிடைக்கும் டியூரிங் பொறியின் நிலை வரைபடம்
  \olref{fig:doubler} இல் காட்டப்பட்டுள்ளது.
  \begin{figure}
\[
\begin{tikzpicture}[->,>=stealth',shorten >=1pt,auto,node distance=2.8cm,
                    semithick]
  \tikzstyle{every state}=[fill=none,draw=black,text=black]

  \node[initial,state] (1)              {$q_0$};
  \node[state]         (2) [right of=1] {$q_1$};
  \node[state]         (3) [right of=2] {$q_2$};
  \node[state]         (4) [below of=3] {$q_3$};
  \node[state]         (5) [left of=4]  {$q_4$};
  \node[state]         (6) [left of=5]  {$q_5$};

  \path (1) edge node {\TMtrans{\TMstroke}{\TMblank}{\TMright}} (2)
    (2) edge [loop above] node {\TMtrans{\TMstroke}{\TMstroke}{\TMright}} (2)
      edge node {\TMtrans{\TMblank}{\TMblank}{\TMright}} (3)
    (3) edge [loop above] node {\TMtrans{\TMstroke}{\TMstroke}{\TMright}} (3)
        edge node {\TMtrans{\TMblank}{\TMstroke}{\TMright}} (4)
    (4) edge [loop below] node {\TMtrans{\TMblank}{\TMstroke}{\TMleft}} (4)
        edge node {\TMtrans{\TMstroke}{\TMstroke}{\TMleft}} (5)
    (5) edge [loop below]  node {\TMtrans{\TMstroke}{\TMstroke}{\TMleft}} (5)
        edge              node {\TMtrans{\TMblank}{\TMblank}{\TMleft}} (6)
    (6) edge [loop below] node {\TMtrans{\TMstroke}{\TMstroke}{\TMleft}} (6)
        edge              node {\TMtrans{\TMblank}{\TMblank}{\TMright}} (1);
\end{tikzpicture}
\]
\caption{ஓர் இரட்டிப்பான் பொறி}
\ollabel{fig:doubler}
\end{figure}
\end{ex}

\begin{prob}
தன்னிச்சையான ஓர் உள்ளீட்டைத் தேர்ந்தெடுத்து,
\olref[tur][mac][rep]{ex:doubler} இலுள்ள இரட்டிப்பான் பொறியின்
அமைவுநிலைகளை ஒன்றன்பின் ஒன்றாகத் தடமறிக.
\end{prob}

\begin{prob}
$\{\TMendtape,\TMblank, A, B\}$ என்ற குறியெழுத்துத் தொகுதியைக்
கொண்ட பின்வரும் டியூரிங் பொறியை வடிவமைக்க: $A$-களின் எண்ணிக்கை
$B$-களின் எண்ணிக்கைக்குச் சமமாகவும் \emph{மேலும்} எல்லா $A$-களும்
எல்லா $B$-களுக்கும் முன் வருவதாகவும் உள்ள எந்த $A$, $B$
எழுத்துத்தொடரையும் பொறி ஏற்க வேண்டும், அதாவது அதில் நிற்க வேண்டும்.
$A$-களின் எண்ணிக்கை~$B$-களின் எண்ணிக்கைக்குச் சமமல்லாத அல்லது
எல்லா~$A$-களும் எல்லா~$B$-களுக்கும் முன் வராத எந்த
எழுத்துத்தொடரையும் பொறி மறுக்க வேண்டும், அதாவது அதில் நிற்கக்கூடாது.
(எடுத்துக்காட்டாக, பொறி $AABB$, $AAABBB$ ஆகியவற்றை ஏற்று,
$AAB$, $AABBAABB$ இரண்டையும் மறுக்க வேண்டும்.)
\end{prob}

\begin{prob}
$\{\TMendtape,\TMblank, A, B\}$ என்ற குறியெழுத்துத் தொகுதியைக்
கொண்ட பின்வரும் டியூரிங் பொறியை வடிவமைக்க: $A$-களும் $B$-களும்
கொண்ட எந்த எழுத்துத்தொடர் $\alpha$ ஐயும் உள்ளீடாகப் பெற்று,
அவற்றை நகலெடுத்து $\alpha\alpha$ என்ற வடிவில் வெளியீட்டை உருவாக்க
வேண்டும். (எடுத்துக்காட்டாக, உள்ளீடு $ABBA$ ஆனது வெளியீடு
$ABBAABBA$ ஐத் தர வேண்டும்.)
\end{prob}

\begin{prob}
\emph{அகரவரிசையிலா?:} $\{\TMendtape,\TMblank, A, B\}$ என்ற
குறியெழுத்துத் தொகுதியைக் கொண்ட பின்வரும் டியூரிங் பொறியை வடிவமைக்க:
$A$-களும் $B$-களும் கொண்ட ஒரு முடிவுறு தொடரை உள்ளீடாகப்
பெறும்போது, எல்லா $A$-களும் எல்லா $B$-களுக்கும் இடப்புறம்
தோன்றுகின்றனவா இல்லையா என்று சரிபார்க்க வேண்டும். பொறி உள்ளீட்டு
எழுத்துத்தொடரை நாடாவில் விட்டுவிட்டு, எழுத்துத்தொடர்
``அகரவரிசையில்'' இருந்தால் நிற்க வேண்டும்; இல்லையெனில்
என்றென்றும் கண்ணியில் இயங்க வேண்டும்.
\end{prob}

\begin{prob}
\emph{அகரவரிசைப்படுத்தி:} $\{\TMendtape,\TMblank, A, B\}$ என்ற
குறியெழுத்துத் தொகுதியைக் கொண்ட பின்வரும் டியூரிங் பொறியை வடிவமைக்க:
$A$-களும் $B$-களும் கொண்ட ஒரு முடிவுறு தொடரை உள்ளீடாகப் பெற்று,
எல்லா $A$-களும் எல்லா~$B$-களுக்கும் இடப்புறம் இருக்குமாறு அவற்றை
மறுவரிசைப்படுத்த வேண்டும். (எடுத்துக்காட்டாக, தொடர் $BABAA$ ஆனது
தொடர் $AAABB$ ஆகவும், தொடர் $ABBABB$ ஆனது தொடர் $AABBBB$ ஆகவும்
மாற வேண்டும்.)
\end{prob}

\end{document}
